\section{Consequences and scope of the theory}

Theorem~\ref{thm:main} is the central result of the paper. It is finite-stage, requires only a finite control-side residual-genomic quantile $b_{\delta,n}$ at the clinical stage being analyzed, and is the part directly instantiated by the Colorado data.

\subsection{Clinical evidence accumulates sequentially}

For a sequential history $C_n=(X_1,\ldots,X_n)$, the clinical likelihood ratio factorizes exactly as
\begin{equation}
\LR_C(C_n)
=\prod_{k=1}^{n}
\frac{p(X_k\mid Y=1,X_{<k})}
     {p(X_k\mid Y=0,X_{<k})},
\label{eq:sequentialLR}
\end{equation}
so that
\begin{equation}
\log\LR_C(C_n)
=\sum_{k=1}^{n}
\log
\frac{p(X_k\mid Y=1,X_{<k})}
     {p(X_k\mid Y=0,X_{<k})}.
\label{eq:sequentiallogLR}
\end{equation}
This is the standard likelihood-ratio accumulation underlying sequential testing and information accumulation \cite{Wald1945,CoverThomas2006}. It shows how additional clinical observations can add evidence. The manuscript makes no asymptotic claim that clinical log-likelihood ratios become mathematically unbounded as history lengthens; the sequential factorization is motivation for longitudinal evidence accumulation, not an additional assumption of Theorem~\ref{thm:main}.

\subsection{Local decision gain need not produce global discrimination gain}

Theorem~\ref{thm:main}(ii) controls where genomics can change a threshold decision, not the global ranking of all case--control pairs. If genomic information matters only in a narrow decision band, a clinically meaningful subset of decisions can change while global AUC changes negligibly. Proposition~\ref{prop:auc-local} formalizes this distinction. This is closely related to the broader biomarker literature showing that strong association need not imply strong discrimination and that incremental AUC can be insensitive when an existing model already performs well \cite{Pepe2004,Pencina2012,Steyerberg2012}. The relevant empirical question is therefore not whether a particular combined learner outperforms ZeBRA globally, but whether genomics contributes reproducible residual evidence after conditioning on the clinical channel and whether that evidence is localized near clinically relevant boundaries.

Weak ZeBRA--genotype association is also not evidence that residual genomic disease information is absent. The relevant object is $\LR_{G\mid C}$, not the marginal dependence between a clinical score and genotype. ZeBRA can be nearly uninformative about rs35705950 carrier status while \textit{MUC5B} still changes disease odds conditional on a given clinical state.

\subsection{Operating thresholds and calibrated posterior scores are special cases}

The decision-relevant genomic region is defined relative to a chosen likelihood-ratio threshold. There is therefore no threshold-free universal ``useful band.'' Corollary~\ref{cor:zebra} translates this geometry into posterior-risk space only when a clinical score is calibrated so that $Z=\mathbb{P}(Y=1\mid C)$. Because raw Colorado ZeBRA is not calibrated as an absolute posterior probability, that corollary is not applied numerically. The empirical analysis instead estimates $\LR_Z$ directly from the observed ZeBRA score. This separation between discrimination and probability calibration follows standard prediction-model evaluation practice \cite{VanCalster2019}.
